Systems built on large language models increasingly adapt by editing the persistent text and code around a frozen model: instructions, skills, memories, tools, and routing. Such systems are often described as self-improving, a label that can cover three different things: an artifact that persists, an edit shown to improve the deployment that made it, and an improvement justified somewhere else. We separate them as two decisions. A local loop keeps an edit only when a check shows that the deployment's failure rate falls; a promotion loop decides how far a kept edit may travel. Reading 256 public systems by who or what makes each decision, we find that the local acceptance check is the most developed piece, while the failure log and the trace from a failure to the responsible artifact are often incomplete; promotion exists only in parts, inside one operator. The feedback a deployment receives (an exact check, an observed outcome, a judgment, or a delayed reward) fixes what a gate can know, and in every class an accepted edit is evidence only for the measure and tasks it was tested on. The reading exposes four confusions the label can hide. Every classification is released row by row.
